Considera les funcions $f(x)=x^2-2x$ i $g(x)=-x^2+4$.

\begin{apartats}

\apartat[1,25]{0,75}
Troba els punts de tall de les dues paràboles, i una primitiva de $g(x)-f(x)$.

\begin{solucio}
$x^2-2x=-x^2+4$, és a dir, $2x^2-2x-4=0$, o $x^2-x-2=0$: $x=-1$ i $x=2$, als punts $(-1,3)$ i $(2,0)$.
$g(x)-f(x)=-2x^2+2x+4$, amb primitiva $-\dfrac{2x^3}3+x^2+4x$.
\end{solucio}

\begin{tria}{area-paraboles}
\itemtria{area}{1}{1,25}
Calcula l'àrea de la regió tancada per les dues paràboles.

\begin{solucio}
Entre $-1$ i 2, $g\ge f$ (a $x=0$, $g=4>f=0$):
\[
\int_{-1}^2(-2x^2+2x+4)\,dx=\left[-\frac{2x^3}3+x^2+4x\right]_{-1}^2=\frac{20}3-\left(-\frac73\right)=9
\]
unitats quadrades.
\end{solucio}

\itemtria{amb-rectes-verticals}{1}{1,25}
Calcula l'àrea de la regió limitada per les gràfiques de $y=x^2$ i $y=x$ i les rectes $x=0$ i $x=2$.

\begin{solucio}
Les dues corbes es tallen a $x=0$ i $x=1$, i la de sobre canvia: a $[0,1]$, $x\ge x^2$, i a $[1,2]$,
$x^2\ge x$.
\[
\int_0^1(x-x^2)\,dx+\int_1^2(x^2-x)\,dx=\frac16+\frac56=1
\]
unitat quadrada.
\end{solucio}
\end{tria}

\begin{nomesllarg}
\apartat{0,75}
Calcula l'àrea de la regió tancada per la gràfica de $f$ i l'eix $X$.

\begin{solucio}
$x^2-2x=0$ dona $x=0$ i $x=2$, i entre aquests valors $f\le0$:
$\left|\displaystyle\int_0^2(x^2-2x)\,dx\right|=\left|\frac83-4\right|=\frac43$ unitats quadrades.
\end{solucio}
\end{nomesllarg}

\end{apartats}
